\subsection{System Model}

We consider a remote IoT monitoring environment comprising a set \(\mathcal{N} = \{1, \ldots, N\}\)
of static sensor devices uniformly distributed over a rectangular Region of Interest (RoI)
\(\mathcal{W} \subset \mathbb{R}^2\). These devices continuously generate monitoring data,
which accumulates in finite transmission buffers. Unlike simple deadline-based models, we
assume data carries value only if it is collected before the buffer capacity \(B_{\max}\) is
exceeded, after which packet loss occurs.

The system model is characterised by a rotary-wing UAV acting as a mobile LoRaWAN gateway,
dispatched to harvest data from the ground sensors. The mission duration is fixed to a finite
horizon \(T\) and divided into discrete time steps \(t = 0, 1, \ldots, T\), each of duration
\(\Delta t\). We assume the UAV flies at a fixed altitude \(h\) to maintain a consistent
ground-to-air channel geometry, updating its horizontal position \(p_t = (x_t, y_t)\) at
each step. Critically, the system employs the LoRaWAN protocol with orthogonal Spreading
Factors (SF), allowing multiple sensors to upload data simultaneously provided they are
assigned distinct SFs by the Adaptive Data Rate (ADR) mechanism. The UAV must therefore
optimise its trajectory not merely to minimise distance, but to actively manage these spectral
assignments to prevent network congestion.

\subsubsection{UAV Kinematics and Energy Consumption}

We consider a rotary-wing Unmanned Aerial Vehicle (UAV) operating within a discretised 3D
grid space \(\mathcal{W} \subset \mathbb{R}^3\). The mission is defined over a finite time
horizon \(T\), discretised into time steps of duration \(\Delta t\). At each time step \(k\),
the UAV's kinematic state is defined by its 3D position \(\mathbf{p}_k = [x_k, y_k, z_k]^T\).

The UAV is controlled by a discrete action space \(\mathcal{A}\) consisting of five high-level
motion primitives:
\begin{equation}
    \mathcal{A} = \{\text{MoveN},\ \text{MoveS},\ \text{MoveE},\ \text{MoveW},\ \text{Hover}\}
    \label{eq:action_space}
\end{equation}
Each movement action updates the UAV's position by a fixed step size \(\delta_d\) in the
corresponding cardinal direction. The Hover action maintains the current position
\(\mathbf{p}_{k+1} = \mathbf{p}_k\), which is critical for stabilising the LoRaWAN link to
facilitate data packet reception.

The energy consumption model is state-dependent. Let \(E_k\) denote the remaining battery
energy at step \(k\). The energy evolution is governed by:
\begin{equation}
    E_{k+1} = E_k - P(a_k) \cdot \Delta t \label{eq:energy_evolution}
\end{equation}
where \(P(a_k)\) is the instantaneous power consumption. We define distinct power profiles
for flight modes:
\begin{equation}
    P(a_k) =
    \begin{cases}
        P_{\text{hover}}  & \text{if } a_k = \text{Hover} \\
        P_{\text{flight}} & \text{if } a_k \in \{\text{Move}\}
    \end{cases}
    \label{eq:power_model}
\end{equation}
Consistent with rotary-wing aerodynamics, we model the power consumption such that
\(P_{\text{hover}} > P_{\text{flight}}\) due to the lack of translational lift in static
flight.\footnote{In rotary-wing UAVs, translational lift---the additional lift gained from
forward flight---means that hovering requires more power than forward flight at moderate
speeds.} Accordingly, we assign a higher power cost to the Hover action (\SI{700}{W})
compared to translational Move actions (\SI{500}{W}). This power differential incentivises
the agent to optimise for efficient flight trajectories and penalises excessive loitering,
thereby aligning the reward signal with the physical reality of limited onboard energy
storage.

\subsubsection{LoRaWAN Channel and Protocol Modelling}

To accurately simulate the ground-to-air communication channel, a simple Free Space Path
Loss (FSPL) model is insufficient as it neglects ground reflections, which are significant at
low UAV altitudes. Instead, this work employs the Two-Ray Ground Reflection Model~\cite{rappaport2002wireless},
which is particularly appropriate for low-altitude UAV platforms where the LoS path and
ground-reflected path interfere constructively or destructively depending on the UAV's
altitude and ground distance~\cite{alhourani2014optimal}.

As shown in \autoref{fig:tworay}, the model calculates the path loss based on the
interference between the direct and reflected rays. The received power \(P_r\) is given by:
\begin{equation}
    P_r = P_t G_t G_r \frac{h_t^2 h_r^2}{d^4} \label{eq:tworay}
\end{equation}
where \(h_t\) and \(h_r\) represent the heights of the transmitter (sensor) and receiver
(UAV), respectively. This results in a steeper signal decay (\(1/d^4\)) compared to free
space, enforcing realistic constraints on the UAV's data collection radius.

\begin{figure}[h]
    \centering
    \fbox{\parbox{0.85\textwidth}{\centering\vspace{1.5cm}
        \textit{[Figure: Two-Ray Ground Reflection Model geometry. Shows the direct LoS path
        and the ground-reflected path between the ground sensor at height $h_t$ and the UAV
        at height $h_r$, illustrating the constructive/destructive interference pattern that
        produces the $1/d^4$ path-loss law beyond the breakpoint distance.]}
    \vspace{1.5cm}}}
    \caption[Two-Ray Ground Reflection Model for UAV-IoT Communication]{Two-Ray Ground
    Reflection Model for UAV-IoT Communication. The model accounts for both the direct
    Line-of-Sight (LoS) path and the ground-reflected path. The destructive interference
    between these two paths beyond the breakpoint distance $d_{\text{break}} = 4\pi h_t h_r / \lambda$
    produces the steeper $1/d^4$ signal decay that enforces realistic data-collection radius
    constraints in the simulation.}
    \label{fig:tworay}
\end{figure}

\subsubsection{Distance and Propagation Geometry}

The simulation operates in a 3D Cartesian coordinate system where the UAV's position at
time \(t\) is denoted as \(\mathbf{p}_{u,t} = (x_u, y_u, h_u)\) and the \(i\)-th sensor's
position is fixed at \(\mathbf{p}_{s,i} = (x_i, y_i, 0)\).

The Euclidean distance \(d_{i,t}\) between the UAV and the \(i\)-th sensor is calculated as
the 3D slant range:
\begin{equation}
    d_{i,t} = \sqrt{(x_u - x_i)^2 + (y_u - y_i)^2 + h_u^2} \label{eq:distance}
\end{equation}
where \(h_u\) is the UAV's altitude. This 3D distance is critical for the Two-Ray Ground
Reflection Model, as the path loss depends heavily on the height differential between the
transmitter and receiver.

In the Two-Ray model, the phase difference between the direct Line-of-Sight (LoS) path and
the ground-reflected path is a function of this distance and the reflection coefficient. For
small grazing angles (typical for UAVs flying at altitude \(h_u\) collecting from ground
sensors), the path loss \(L\) simplifies to the fourth-power law:
\begin{equation}
    L \propto \frac{h_u^2 h_s^2}{d_{i,t}^4} \label{eq:fourthpower}
\end{equation}
This quartic decay (\(1/d^4\)) means that small changes in the UAV's position or altitude
have a magnified impact on the Received Signal Strength Indicator (RSSI), creating a highly
sensitive environment that the DRL agent must learn to navigate precisely.

\subsubsection{LoRaWAN Device Class and Traffic Model}

We model the IoT sensor nodes as LoRaWAN Class A devices~\cite{augustin2016lora,bor2016lora},
operating in an asynchronous, sensor-initiated transmission mode. Class A operation imposes a strict energy constraint
where nodes remain in a deep sleep state and activate solely to transmit data packets
according to a stochastic duty cycle. Consequently, the network operates without centralised
synchronisation or wake-up receiver infrastructure.

This design choice reflects the practical reality of energy-harvesting IoT deployments. Class
A devices are the de facto standard in real-world LoRaWAN networks because their simplicity
and low power consumption make them ideal for battery-powered and energy-harvesting
applications. The lack of downlink capability in Class A devices precludes the use of wake-up
receiver architectures, which would require active, power-consuming receivers to remain
listening for network-initiated wake-up signals---a fundamental contradiction with the strict
energy budgets of sensor nodes. Thus, our modelling choice ensures that the proposed
framework remains applicable to genuine large-scale deployments, rather than idealised
systems with unlimited infrastructure support.

\paragraph{Data Generation and Buffer Model.}
Data generation at each sensor \(i \in \mathcal{N}\) follows a Poisson process with rate
\(\lambda_i\):
\begin{equation}
    P(\text{packet generated at sensor } i \text{ in interval } [t, t + \Delta t]) = \lambda_i \cdot \Delta t
    \label{eq:poisson}
\end{equation}
Generated packets accumulate in finite transmission buffers with capacity \(B_{\max}\):
\begin{equation}
    \text{Buffer}_i(t) \in [0,\, B_{\max}] \label{eq:buffer}
\end{equation}
Once a buffer reaches capacity, newly generated packets are discarded without transmission.
This models the realistic scenario wherein sensor data loses value if not collected
promptly---a behaviour common in time-sensitive monitoring applications such as environmental
sensing, disaster response, and industrial monitoring. The buffer overflow mechanism creates a
hard deadline for data collection: if the UAV does not visit a sensor before its buffer fills,
the oldest or newest packets (depending on buffer policy) are lost permanently.

\paragraph{Mission Implications.}
The combination of Class A asynchronous operation and finite buffer capacity creates a
challenging multi-objective optimisation problem for the DRL agent:
\begin{itemize}
    \item \emph{Fairness constraint}: The agent must visit all sensors periodically to prevent
    buffer overflow and information loss.
    \item \emph{Energy efficiency constraint}: The agent must minimise flight distance and
    hovering time to conserve battery energy.
    \item \emph{Throughput objective}: The agent must maximise the total number of packets
    collected within the mission horizon.
\end{itemize}
These competing objectives naturally emerge from the Class A device constraints and form the
basis of the reward function design discussed in \autoref{sec:reward}.

\subsubsection{Collision Resolution via Capture Effect} \label{sec:capture}

While orthogonal spreading factors do not interfere with one another, intra-SF collisions
occur when multiple sensors attempt transmission on the same spreading factor and frequency
simultaneously. We resolve these collisions using the Capture Effect model~\cite{croce2020lora}, which
exploits the power-law properties of LoRa receivers to successfully decode the strongest
signal whilst suppressing weaker interferers.

\paragraph{Capture Effect Criterion.}
The gateway successfully decodes a signal from sensor \(i\) if its received power \(P_{r,i}\)
exceeds the sum of all interfering signals by a capture threshold \(\tau_{\text{cap}}\)
(typically \SI{6}{dB}):
\begin{equation}
    \frac{P_{r,i}}{\displaystyle\sum_{j \in \mathcal{I}_{\text{co-channel}}} P_{r,j} + N_0} \geq \tau_{\text{cap}}
    \label{eq:capture}
\end{equation}
where \(N_0\) is the background noise power spectral density and
\(\mathcal{I}_{\text{co-channel}} = \{j \in \mathcal{N} \mid j \neq i,\ \text{SF}_j(t) = \text{SF}_i(t)\}\)
denotes the set of sensors transmitting on the same spreading factor as sensor \(i\).

\paragraph{Implication for Trajectory Optimisation.}
This constraint is critical to the DRL agent's learning objective. The Capture Effect forces
the agent to optimise its trajectory not solely for coverage---i.e., achieving minimum
received signal strength indicator (RSSI) thresholds---but rather to induce a power
differential that ensures successful packet capture during congestion events.

Consider an intra-SF collision scenario in which two sensors simultaneously transmit on the
same spreading factor. The received powers at the UAV are:
\begin{align}
    P_{r,i} &= P_t G_t G_r \frac{h_t^2 h_r^2}{d_{i,t}^4} \label{eq:pri} \\
    P_{r,j} &= P_t G_t G_r \frac{h_t^2 h_r^2}{d_{j,t}^4} \label{eq:prj}
\end{align}
For sensor \(i\) to be successfully captured, the UAV must position itself such that
\(d_{i,t} < d_{j,t}\), thereby satisfying:
\begin{equation}
    \frac{d_{j,t}^4}{d_{i,t}^4} \geq \tau_{\text{cap}} \label{eq:capture_cond}
\end{equation}
This creates a non-trivial spatial optimisation problem wherein naive distance-minimisation or
greedy path-planning strategies fundamentally fail. The agent must learn to position itself
strategically closer to the desired sensor during collision events, not merely visit all
sensors with equal frequency. This insight is fundamental to the work: the UAV's physical
position becomes an active control variable for the communication protocol, rather than a
passive design parameter.

\subsubsection{Gateway Architecture}

The UAV communication module is modelled as a single-channel multi-spreading-factor gateway.
We exploit the orthogonality of LoRa spreading factors to demodulate multiple concurrent
signals transmitted simultaneously on distinct spreading factors. This property---unique to
LoRa modulation---allows us to achieve concurrent channel capacity beyond what would be
possible in traditional narrowband systems.

\paragraph{Spreading Factor Selection.}
Rather than implementing the full LoRaWAN standard set, which defines spreading factors
\(\text{SF} \in \{7, 8, 9, 10, 11, 12\}\), we restrict the operational spreading factor set
to:
\begin{equation}
    \mathcal{S} = \{7, 9, 11, 12\} \label{eq:sf_set}
\end{equation}
This pragmatic restriction reflects the hardware limitations of real UAV-mounted gateways,
which typically operate a single transceiver with time-division or frequency-division
multiplexing across a restricted subset of spreading factors. The specific selection of
\(\mathcal{S} = \{7, 9, 11, 12\}\) maintains sufficient spectral diversity to prevent the
spreading factor monoculture whilst remaining implementable on lightweight airborne platforms
with constrained power and computational budgets.

The four spreading factors in \(\mathcal{S}\) are mutually orthogonal, meaning signals
modulated with different spreading factors can theoretically coexist on the same frequency
channel without interfering with one another, provided they do not collide on the same
spreading factor and frequency pair.

\paragraph{Concurrent Reception and Collision Handling.}
This single-channel multi-SF architecture enables the UAV to receive up to four simultaneous
packets, subject to the constraint:
\begin{equation}
    |\{i \in \mathcal{N} : \text{SF}_i(t) = k\}| \leq 1, \quad \forall k \in \mathcal{S}
    \label{eq:concurrent}
\end{equation}
where \(\text{SF}_i(t)\) denotes the spreading factor assigned to sensor \(i\) at time \(t\),
and \(\mathcal{N} = \{1, \ldots, N\}\) is the set of all deployed sensor nodes.

\autoref{eq:concurrent} ensures that at most one sensor transmits on each spreading factor at
any given time. Packets that violate this constraint are subject to collision resolution via
the Capture Effect mechanism detailed in \autoref{sec:capture}.
